\documentclass[11pt]{article}
\usepackage[margin=1in]{geometry}
\usepackage{amsmath,amssymb,amsthm}
\usepackage{booktabs,array,float}
\newtheorem{theorem}{Theorem}[section]

\title{Increasing Logarithmic Weights for Rolling Validation: Theory}
\author{Aileen Tao}
\date{}

\begin{document}
\maketitle

\section{Introduction and research question}

Ordinary rolling validation adds validation losses across successive
validation times.  Here later observations receive larger weights.  The
question is how this changes the expected advantage of one model, the
variability of that comparison, and the probability of selecting the wrong
model.  In particular, we ask which forms of increasing logarithmic growth
still allow weighted rolling validation to favor the oracle consistently.
The companion report, \emph{Increasing Logarithmic Weights for Rolling
Validation: Simulation Study} (08B), examines the finite-sample behavior.

\subsection{Statistical setting}

Let $Y_i=\theta_\star+\xi_i$, where
$\xi_i\stackrel{\mathrm{iid}}{\sim}N(0,1)$.  At validation time $i$, compare
Model 1, the oracle $\widehat\theta_{i-1}^{(1)}=\theta_\star$, with Model 0,
the fitted mean
\[
\widehat\theta_{i-1}^{(0)}=\frac{1}{i-1}\sum_{t=1}^{i-1}Y_t.
\]
Model 1 uses the true parameter, while Model 0 estimates it from the
observations available before time $i$.  This comparison isolates the
finite-sample cost of estimation in a simple setting.

\subsection{Weighted rolling-validation statistic and selection criterion}

For positive deterministic weights $w_i$, define
\begin{equation}
\Xi_n=\sum_{i=2}^n w_i\left[
\bigl(Y_i-\widehat\theta_{i-1}^{(0)}\bigr)^2
-\bigl(Y_i-\widehat\theta_{i-1}^{(1)}\bigr)^2\right].
\label{eq:statistic}
\end{equation}
A positive $\Xi_n$ means Model 0 has the larger weighted loss, so the oracle
is selected.  A negative $\Xi_n$ favors Model 0 and is therefore an incorrect
selection.  Selection consistency means $\mathbb P(\Xi_n<0)\to0$.

We study each weight through three quantities: $\mathbb E[\Xi_n]$, the
accumulated signal favoring the oracle; $\operatorname{Var}(\Xi_n)$, the
variability of that comparison; and their ratio
$\operatorname{Var}(\Xi_n)/\{\mathbb E[\Xi_n]\}^2$.

\paragraph{Expected signal.}
For any positive deterministic weights,
\begin{equation}
\mathbb E[\Xi_n]=\sum_{i=2}^n\frac{w_i}{i-1}.
\label{eq:mean}
\end{equation}
This is the accumulated deterministic advantage of the oracle under the
chosen weighting sequence.

\paragraph{Variability.}
Successive validation terms share training observations, so their dependence
must be retained.  The exact variance decomposition and its Gaussian
orthogonality argument are given in Appendix~A.

\paragraph{Chebyshev selection rule.}
Since $\mathbb E[\Xi_n]>0$,
\begin{equation}
\mathbb P(\Xi_n<0)
\leq\frac{\operatorname{Var}(\Xi_n)}{\{\mathbb E[\Xi_n]\}^2}.
\label{eq:chebyshev}
\end{equation}
Thus the common three-step calculation is to find the mean, find the
variance, and test whether their Chebyshev ratio tends to zero.  A vanishing
ratio is sufficient for the wrong-selection probability to vanish.

\subsection{Candidate weights}

We compare four increasing weights with $w_i=1$, which serves only as an
unweighted baseline and is not part of the increasing-weight family.

\begin{table}[H]
\centering
\begin{tabular}{@{}ll@{}}
\toprule
Weight & Why it is included \\
\midrule
$1$ & Unweighted baseline \\
$\log(i+1)$ & Slowly increasing logarithmic weight \\
$[\log(i+1)]^2$ & Stronger logarithmic growth \\
$1+\log(i+1)+[\log(i+1)]^2$ & Polynomial-log case \\
$\log\log(i+e)$ & Extremely slowly increasing case \\
\bottomrule
\end{tabular}
\caption{Weights considered in 08A and 08B, in their common order.}
\end{table}

\section{Positive log-power weights}

The family $w_i=[\log(i+1)]^p$, $p>0$, provides one framework for different
logarithmic growth rates.  The exponent controls the relative emphasis on
later validation observations.  The question is whether the extra
variability from increasing weights is dominated by the accumulated signal.

\begin{theorem}\label{thm:logpower}
For every fixed $p>0$,
\[
\textnormal{Mean:}\qquad
\mathbb E[\Xi_n]\sim\frac{(\log n)^{p+1}}{p+1},
\]
\[
\textnormal{Variance:}\qquad
\operatorname{Var}(\Xi_n)\sim2(\log n)^{2p},
\]
and
\[
\textnormal{Selection:}\qquad
\frac{\operatorname{Var}(\Xi_n)}{\{\mathbb E[\Xi_n]\}^2}
\sim\frac{2(p+1)^2}{(\log n)^2}\longrightarrow0.
\]
Hence every fixed $p>0$ gives selection consistency.
\end{theorem}

For every fixed $p>0$, the ratio has order $(\log n)^{-2}$ and tends to zero.
Weighting increases both signal and variance, but the squared mean grows two
powers of $\log n$ faster.  The full proof is given in Appendix~B.

\subsection{Two concrete choices}

The log and squared-log weights are direct special cases.
\begin{center}
\begin{tabular}{@{}lcccl@{}}
\toprule
Weight & $\mathbb E[\Xi_n]$ & $\operatorname{Var}(\Xi_n)$
& Chebyshev ratio & Selection \\
\midrule
$\log(i+1)$ & $\sim(\log n)^2/2$ & $\sim2(\log n)^2$
& $\sim8/(\log n)^2$ & Consistent \\
$[\log(i+1)]^2$ & $\sim(\log n)^3/3$ & $\sim2(\log n)^4$
& $\sim18/(\log n)^2$ & Consistent \\
\bottomrule
\end{tabular}
\end{center}
Their finite-sample moments, ratios, and selection errors are examined in
08B, Sections 3--5.

\section{Polynomial-log weights}

The positive log-power result suggests that the dominant growth rate may
determine the asymptotic behavior.  This motivates a polynomial in
$\log(i+1)$.

\begin{theorem}\label{thm:poly}
Suppose $w_i=P(\log(i+1))>0$, where
$P(x)=cx^k+O(x^{k-1})$, $k\geq1$, and $c>0$.  Then
\[
\textnormal{Mean:}\qquad
\mathbb E[\Xi_n]\sim\frac{c}{k+1}(\log n)^{k+1},
\]
\[
\textnormal{Variance:}\qquad
\operatorname{Var}(\Xi_n)\sim2c^2(\log n)^{2k},
\]
and
\[
\textnormal{Selection:}\qquad
\frac{\operatorname{Var}(\Xi_n)}{\{\mathbb E[\Xi_n]\}^2}
\sim\frac{2(k+1)^2}{(\log n)^2}\longrightarrow0.
\]
\end{theorem}

The lower-degree terms grow more slowly than the leading $c[\log(i+1)]^k$
term, so they do not change the leading mean or variance.  In particular,
$1+\log(i+1)+[\log(i+1)]^2$ has the same leading asymptotics as
$[\log(i+1)]^2$.  See Appendix~C for the proof and 08B, Sections 3--5 for
the finite-sample comparison.

\section{The log-log weight}

The weight $w_i=\log\log(i+e)$ tests whether an even more slowly increasing
sequence can still give a vanishing Chebyshev ratio.

\begin{theorem}\label{thm:loglog}
For $w_i=\log\log(i+e)$,
\[
\textnormal{Mean:}\qquad
\mathbb E[\Xi_n]\sim(\log n)\log\log n,
\]
\[
\textnormal{Variance:}\qquad
\operatorname{Var}(\Xi_n)\sim2(\log\log n)^2,
\]
and
\[
\textnormal{Selection:}\qquad
\frac{\operatorname{Var}(\Xi_n)}{\{\mathbb E[\Xi_n]\}^2}
\sim\frac{2}{(\log n)^2}\longrightarrow0.
\]
\end{theorem}

Even this extremely slow increase is sufficient in the present example.
The proof is in Appendix~D, and the matching finite-sample evidence is in
08B, Sections 3--5.

\section{Theory summary}

\begin{table}[H]
\centering
\footnotesize
\setlength{\tabcolsep}{3pt}
\begin{tabular}{@{}>{\raggedright\arraybackslash}p{3.0cm}
>{\raggedright\arraybackslash}p{2.8cm}
>{\raggedright\arraybackslash}p{2.8cm}
>{\raggedright\arraybackslash}p{3.0cm}
>{\raggedright\arraybackslash}p{1.7cm}@{}}
\toprule
Weight & $\mathbb E[\Xi_n]$ & $\operatorname{Var}(\Xi_n)$
& Chebyshev ratio & Selection \\
\midrule
$1$ & $\sim\log n$ & $\to\pi^2$ & $\sim\pi^2/(\log n)^2$ & Consistent \\
$\log(i+1)$ & $\sim(\log n)^2/2$ & $\sim2(\log n)^2$
& $\sim8/(\log n)^2$ & Consistent \\
$[\log(i+1)]^2$ & $\sim(\log n)^3/3$ & $\sim2(\log n)^4$
& $\sim18/(\log n)^2$ & Consistent \\
$1+\log(i+1)+[\log(i+1)]^2$ & $\sim(\log n)^3/3$
& $\sim2(\log n)^4$ & $\sim18/(\log n)^2$ & Consistent \\
$\log\log(i+e)$ & $\sim(\log n)\log\log n$
& $\sim2(\log\log n)^2$ & $\sim2/(\log n)^2$ & Consistent \\
\bottomrule
\end{tabular}
\caption{Theoretical predictions examined in the accompanying simulation report.}
\end{table}

\section{Conclusion}

Every fixed positive log-power weight considered here remains selection
consistent because its accumulated signal dominates its variability.  For
polynomial-log weights, the highest-degree term determines the leading
behavior.  The log-log result shows that even extremely slow increasing
growth can be sufficient.

The signal must grow sufficiently fast relative to its standard deviation:
\[
\frac{\sqrt{\operatorname{Var}(\Xi_n)}}{\mathbb E[\Xi_n]}\longrightarrow0,
\qquad\text{equivalently}\qquad
\frac{\operatorname{Var}(\Xi_n)}{\{\mathbb E[\Xi_n]\}^2}\longrightarrow0.
\]
This signal-versus-noise condition is what makes the Chebyshev consistency
argument work.  The accompanying 08B report checks how clearly these
asymptotic predictions appear at finite sample sizes.

\clearpage
\appendix
\section*{Appendix}

\section{General mean and variance decomposition}

Define $S_m=\sum_{t=1}^m\xi_t$.  Then
\begin{equation}
\Xi_n=\sum_{i=2}^n w_i\left\{
\frac{S_{i-1}^2}{(i-1)^2}-2\xi_i\frac{S_{i-1}}{i-1}\right\}.
\label{eq:expanded}
\end{equation}
Because $S_{i-1}$ is independent of $\xi_i$,
$\mathbb E(\xi_iS_{i-1})=0$ and $\mathbb E(S_{i-1}^2)=i-1$, proving
\eqref{eq:mean}.

For $1\leq t<n$ and $2\leq l\leq n$, define
\[
a_{n,t}=\sum_{i=t+1}^n\frac{w_i}{(i-1)^2},\qquad
b_{n,l}=\sum_{i=l+1}^n\frac{w_i}{(i-1)^2}-\frac{w_l}{l-1}.
\]
Expanding $S_{i-1}^2$ and collecting Gaussian monomials gives
\begin{equation}
\Xi_n-\mathbb E[\Xi_n]
=\sum_{t=1}^{n-1}a_{n,t}(\xi_t^2-1)
+2\sum_{l=2}^n b_{n,l}\sum_{k<l}\xi_k\xi_l.
\label{eq:decomp}
\end{equation}
The centered squares have variance $2$, and each product $\xi_k\xi_l$ has
variance $1$.  Distinct terms are uncorrelated because some independent
centered Gaussian variable appears to an odd power in every mixed product.
There are $l-1$ products with larger index $l$.  Hence
\begin{equation}
\operatorname{Var}(\Xi_n)=2\sum_{t=1}^{n-1}a_{n,t}^2
+4\sum_{l=2}^n(l-1)b_{n,l}^2.
\label{eq:variance}
\end{equation}

For the increasing weights below, define
\[
A_t=\sum_{i=t+1}^\infty\frac{w_i}{(i-1)^2},\qquad
B_l=A_l-\frac{w_l}{l-1},\qquad T_n=A_n.
\]
Then $a_{n,t}=A_t-T_n$ and $b_{n,l}=B_l-T_n$ exactly.  The leading terms
in $B_l$ cancel.

\section{Proof for increasing log-power weights}

\begin{proof}[Proof of Theorem~\ref{thm:logpower}]
Replacing $[\log(i+1)]^p/(i-1)$ by $(\log i)^p/i$ changes the mean by a
bounded amount.  With $u=\log x$, integral comparison gives
\[
\mathbb E[\Xi_n]\sim\int^n\frac{(\log x)^p}{x}\,dx
=\frac{(\log n)^{p+1}}{p+1}.
\]

For $I_p(t)=\int_t^\infty(\log x)^p x^{-2}\,dx$, integration by parts gives
$I_p(t)=(\log t)^p/t+pI_{p-1}(t)$.  Substituting $x=te^v$ shows
$I_q(t)=O((\log t)^q/t)$ for fixed $q$.  Sum--integral comparison therefore
yields
\begin{equation}
A_t\sim\frac{(\log t)^p}{t},\qquad
T_n\sim\frac{(\log n)^p}{n}.
\label{eq:tails}
\end{equation}

To expose the cancellation in $B_l$, replace $(i-1)^{-2}$ in $A_l$ by
$[i(i-1)]^{-1}$; the error is $O((\log l)^p/l^2)$.  Summation by parts gives
\[
\sum_{i=l+1}^\infty\frac{w_i}{i(i-1)}
=\frac{w_{l+1}}{l}+\sum_{i=l+2}^\infty\frac{w_i-w_{i-1}}{i-1}.
\]
The first term differs from $w_l/(l-1)$ by
$O((\log l)^p/l^2+(\log l)^{p-1}/l^2)$, and the mean-value theorem gives
$|w_i-w_{i-1}|\leq C_p(\log i)^{p-1}/i$.  Thus
\begin{equation}
|B_l|\leq C_p\left\{\frac{(\log l)^{p-1}}{l}
+\frac{(\log l)^p}{l^2}\right\}.
\label{eq:Bbound}
\end{equation}
In particular, $\sum_tA_t^2<\infty$, so the diagonal part of
\eqref{eq:variance} is bounded.

For the nonterminal off-diagonal part,
\[
\sum_{l=2}^n(l-1)B_l^2
=O\left(\sum_{l=2}^n\frac{(\log l)^{2p-2}}{l}\right).
\]
With $u=\log x$, the comparison integral is
$\int^{\log n}u^{2p-2}\,du$.  It is bounded for $0<p<1/2$, of order
$\log\log n$ for $p=1/2$, and of order $(\log n)^{2p-1}$ for $p>1/2$.
In every case,
\begin{equation}
\sum_{l=2}^n(l-1)B_l^2=o((\log n)^{2p}).
\label{eq:Blower}
\end{equation}

The terminal contribution satisfies
\[
T_n^2\sum_{l=2}^n(l-1)\sim\frac{(\log n)^{2p}}2.
\]
After expanding $(B_l-T_n)^2$, Cauchy--Schwarz bounds the cross term by
\begin{align*}
2\left|T_n\sum_{l=2}^n(l-1)B_l\right|
&\leq2\left\{
\left(\sum_{l=2}^n(l-1)B_l^2\right)
\left(T_n^2\sum_{l=2}^n(l-1)\right)\right\}^{1/2}\\
&=o((\log n)^{2p}).
\end{align*}
The bounded diagonal part, the $B_l^2$ part, and the cross term are lower
order.  Therefore
\[
\operatorname{Var}(\Xi_n)
\sim4T_n^2\sum_{l=2}^n(l-1)\sim2(\log n)^{2p}.
\]
Dividing by the squared mean completes the proof.
\end{proof}

\section{Polynomial-log extension}
\begin{proof}[Proof of Theorem~\ref{thm:poly}]
Since $P(x)=cx^k+O(x^{k-1})$,
\[
P(\log(i+1))=c[\log(i+1)]^k+O([\log(i+1)]^{k-1}).
\]
Integral comparison gives
$\mathbb E[\Xi_n]=c(\log n)^{k+1}/(k+1)+O((\log n)^k)$.
The tail argument above yields
\[
A_t\sim\frac{c(\log t)^k}{t},\qquad
B_l=O\left(\frac{(\log l)^{k-1}}{l}+\frac{(\log l)^k}{l^2}\right).
\]
Hence $T_n\sim c(\log n)^k/n$,
$\sum_{l=2}^n(l-1)B_l^2=O((\log n)^{2k-1})$, and the cross term is lower
order by Cauchy--Schwarz.  The bounded diagonal term is also lower order, so
\[
\operatorname{Var}(\Xi_n)
\sim4T_n^2\sum_{l=2}^n(l-1)\sim2c^2(\log n)^{2k}.
\]
\end{proof}

\section{Log-log extension}
\begin{proof}[Proof of Theorem~\ref{thm:loglog}]
Integral comparison and $u=\log x$ give
\[
\mathbb E[\Xi_n]
=\int^n\frac{\log\log x}{x}\,dx+O(1)
=(\log n)\log\log n-\log n+O(1).
\]
The same tail comparison and summation by parts give
\[
A_t\sim\frac{\log\log t}{t},\qquad
B_l=O\left(\frac1{l\log l}+\frac{\log\log l}{l^2}\right).
\]
Both $\sum_tA_t^2$ and $\sum_l(l-1)B_l^2$ converge, while
$T_n\sim\log\log n/n$.  The cross term is lower order by Cauchy--Schwarz.
Thus
\[
\operatorname{Var}(\Xi_n)
\sim4T_n^2\sum_{l=2}^n(l-1)\sim2(\log\log n)^2,
\]
which proves the result.
\end{proof}

\end{document}
